\documentclass[11pt]{amsart}
\usepackage{amsmath, amssymb}
\usepackage{hyperref}
\begin{document}

The ``hence'' is invoking a standard principle --- that the Taylor coefficients of a
function (the coefficients of its expansion about $0$) are controlled by its singularity
nearest the origin (the \emph{dominant singularity})~\cite{Hadamard1892,FlajoletSedgewick2009}.
Let me unpack why that principle applies and what it gives, because it is really two claims
stacked together, one elementary and rigorous, the other standard but leading-order.

\subsection*{Setup}
The $g_m$ are the Taylor coefficients of $G$ about the origin: $G(t)=\sum_{m\ge0} g_m t^m$.
The preceding sentence established that $G$'s singularities (the places where this series
stops representing $G$) are the zeros of $H$, the nearest being $t_0$ at distance
$\rho=|t_0|\approx0.043$.

\subsection*{Which Taylor expansion ``the'' refers to}
By ``the Taylor coefficients'' I mean the coefficients of the expansion about $0$ --- the
$g_m$ in $G(t)=\sum_{m\ge0}g_m t^m$. A function has a different Taylor expansion about every
point at which it is analytic, so the phrase is meaningful only once a center is fixed; here
the center is $0$, and I should have said so. Three reasons make $0$ the center that
matters, in increasing order of importance.

\emph{(i) It is the expansion we are after.} $G$ is defined as a power series in $t$ about
the origin: the $g_m=h_m^{(0)}$ are produced by the recurrence (D) from the seed $g_0=1$,
which is exactly the expansion at $0$. The deformed coefficients
$h_n^{(c)}=\sum_k\frac{c^k}{k!}g_{n-k}$, the characteristic polynomials, and the minimum
moduli are all built from these $g_m$. No other expansion enters the problem, so ``the
Taylor coefficients'' can only mean these.

\emph{(ii) It is forced by the normalization.} The sequences are pinned at $h_0=1$, i.e.\
$G(0)=1$; the generating function is by construction organized about $t=0$. Expanding about
another point would not produce the sequence under study.

\emph{(iii) The radius --- and hence which singularity is ``nearest'' --- is defined
relative to the center.} The expansion about $0$ converges in the largest disc
\emph{centered at $0$} containing no singularity; its radius is the distance from $0$ to the
closest zero of $H$, namely $\rho=|t_0|$~\cite{Ahlfors,Hadamard1892}. This is what makes
$t_0$ the nearest singularity, and what ties $|g_m|^{1/m}\to1/\rho$ to $t_0$ specifically.
The phrase ``nearest the origin'' in the main document already encodes the choice of center:
nearest \emph{to $0$}. Expanding about some other point $a$ would recenter the disc at $a$,
make ``nearest'' mean nearest to $a$, change the radius to the distance from $a$ to the
closest zero, and govern those (different) coefficients by a possibly different zero.

So the unsloppy phrasing is \emph{the coefficients of the Taylor expansion of $G$ about
$0$}, and it is the one that matters because it is the very sequence $g_m$ the problem is
made of --- and the growth law $|g_m|^{1/m}\to1/\rho$ is meaningful only once the center
(here $0$) fixes what ``radius'' and ``nearest singularity'' refer to.

\subsection*{Layer 1 --- the growth rate (rigorous, elementary)}
A power series converges exactly up to the nearest singularity of the function it
represents; its radius of convergence equals the distance to that
singularity~\cite{Ahlfors,Hadamard1892}. So the radius of $\sum g_m t^m$ is $\rho$. By the
Cauchy--Hadamard formula (radius $=1/\limsup_m|g_m|^{1/m}$)~\cite{Ahlfors}, this is the same
as saying
\[
|g_m|^{1/m}\longrightarrow \frac1\rho,\qquad\text{i.e.}\qquad |g_m|\ \text{grows like}\ \rho^{-m}.
\]
That much is airtight, and it is exactly what the growth-rate figure of the companion note
confirms ($|g_m|^{1/m}\to 23\approx1/\rho$). So the \emph{modulus} of the nearest zero of
$H$ sets the exponential growth rate. Already this is one sense of ``governed by the nearest
zero.''

\subsection*{Layer 2 --- why the nearest one, and the phase (standard, leading-order)}
Why does the \emph{nearest} singularity dominate, rather than all of them mattering equally?
Because any singularity at a point $t_k$ contributes a term of size $\sim t_k^{-m}$ to the
$m$-th coefficient~\cite{FlajoletOdlyzko1990,FlajoletSedgewick2009}. (A simple example: a
pole, $\frac{1}{t_k-t}=\frac{1}{t_k}\sum_m (t/t_k)^m$, has coefficients $t_k^{-m-1}$; a
branch point $(1-t/t_k)^{\alpha}$ has coefficients $\sim t_k^{-m}m^{-\alpha-1}$. In both
cases the geometric factor is $t_k^{-m}$.) So heuristically
\[
g_m\;\approx\;\sum_k c_k\, t_k^{-m}\times(\text{slowly varying}),
\]
and the term with the \emph{smallest} $|t_k|$ wins, because $|t_k^{-m}|=\rho_k^{-m}$ is
largest when $\rho_k$ is smallest. The next zero sits at $\rho_2\approx0.078$, so its
contribution is smaller than the nearest by a factor
$(\rho/\rho_2)^m=(0.043/0.078)^m\approx 0.56^m$ --- already negligible by $m\approx20$. That
exponential suppression of everything but the nearest singularity is the content of
``hence.''

Now the \emph{phase}. Near a simple zero $t_0$ of $H$, a short calculation gives
$J=tH'/H\approx t_0/(t-t_0)$~\cite{Ahlfors}, so $\int J\approx t_0\log(t-t_0)$ and
$G\approx(t-t_0)^{t_0}$ --- a branch point. Its coefficient contribution carries the
geometric factor $t_0^{-m}$, and since $t_0=\rho e^{i\theta}$,
\[
t_0^{-m}=\rho^{-m}e^{-im\theta}.
\]
The modulus $\rho^{-m}$ is the growth (Layer~1); the unit-modulus part $e^{-im\theta}$ is a
rotation by angle $\theta$ each step. Because $H$ has real coefficients its zeros come in
conjugate pairs $t_0,\bar t_0$, and adding the two conjugate contributions turns
$e^{\pm im\theta}$ into a real cosine:
\[
g_m\;\sim\;C\,\rho^{-m}\cos(m\theta+\varphi).
\]
So the \emph{argument} $\theta$ of the nearest zero sets the
oscillation~\cite{Odlyzko1982,FlajoletSedgewick2009}. That is the second, and for us the
important, sense of ``governed by the nearest zero.''

\subsection*{So the ``hence'' says}
Of all the branch points of $G$ (= zeros of $H$), only the nearest matters for large $m$,
because the others are exponentially smaller; and that nearest one delivers both pieces of
the asymptotics --- its modulus $\rho$ the growth rate $\rho^{-m}$, its argument $\theta$ the
oscillation $\cos(m\theta+\varphi)$. Everything downstream (the period-$4$ in $\mu_n$, the
$\approx37$ envelope) is read off $\theta$.

\subsection*{Honest qualification}
Layer~1 is rigorous. Layer~2 --- that the nearest singularity's \emph{local type} fixes the
full leading form including the polynomial factor and phase --- is the transfer principle of
singularity analysis (Flajolet--Odlyzko)~\cite{FlajoletOdlyzko1990,FlajoletSedgewick2009},
which is a theorem under regularity hypotheses: the singularity should be isolated and of
tractable (algebraic/logarithmic) type, with $G$ continuable in a neighborhood beyond the
disc (a $\Delta$-domain). Those hold plausibly here --- $t_0$ is isolated, well inside the
unit disc, and a genuine branch point --- and the prediction is confirmed numerically, but I
have not verified the technical hypotheses for this particular $G$ (in part because the
behavior of $H$ out on $|t|=1$, where the $\tau(p_n)$ live, is not something I control). So
treat Layer~2 as ``the standard principle, leading order, empirically confirmed,'' not as a
proved theorem.

\begin{thebibliography}{9}
\bibitem{Ahlfors} L.~V.~Ahlfors, \emph{Complex Analysis}, 3rd ed., McGraw--Hill, New York, 1979.
  (Radius of convergence as distance to the nearest singularity; the Cauchy--Hadamard
  formula; the logarithmic derivative and its residue at a zero.)
\bibitem{Hadamard1892} J.~Hadamard, \emph{Essai sur l'\'etude des fonctions donn\'ees par
  leur d\'eveloppement de Taylor}, J. Math. Pures Appl. (4) \textbf{8} (1892), 101--186.
  (Foundational study of how Taylor coefficients reflect singularities.)
\bibitem{Odlyzko1982} A.~M.~Odlyzko, \emph{Periodic oscillations of coefficients of power
  series that satisfy functional equations}, Adv. in Math. \textbf{44} (1982), 180--205.
\bibitem{FlajoletOdlyzko1990} P.~Flajolet and A.~M.~Odlyzko, \emph{Singularity analysis of
  generating functions}, SIAM J. Discrete Math. \textbf{3} (1990), no.~2, 216--240.
\bibitem{FlajoletSedgewick2009} P.~Flajolet and R.~Sedgewick, \emph{Analytic Combinatorics},
  Cambridge University Press, Cambridge, 2009. (Dominant singularities, Ch.~IV--V;
  singularity analysis and transfer theorems, Ch.~VI.)
\end{thebibliography}

\end{document}
